\documentclass[10pt,a4paper]{article}
\usepackage[utf8]{inputenc}
\usepackage[T1]{fontenc}
\usepackage[french]{babel}
\usepackage[margin=1.2cm]{geometry}
\usepackage{enumitem}
\setlength{\parindent}{0pt}
\pagestyle{empty}
\newcommand{\cvsection}[1]{\par\vspace{5pt}\textbf{#1}\\[-2pt]\rule{\linewidth}{0.3pt}\vspace{3pt}}
\begin{document}
\begin{center}
{\Large \textbf{CAMILLE MARTIN}}\\[3pt]
\textbf{Ingénieure systèmes | Python, données | Lyon}\\[3pt]
Email : camille.martin@example.test | Téléphone : +33 1 00 00 00 01
\end{center}
\cvsection{Résumé}
Ingénieure spécialisée dans les systèmes de données fiables et l'automatisation.
\cvsection{Expérience}
\textbf{GridLab -- Ingénieure données} \hfill 2024--2026
\begin{itemize}[leftmargin=*]
\item Développement de pipelines Python et SQL pour des données énergétiques.
\item Contrôles qualité et supervision des traitements.
\end{itemize}
\cvsection{Projets}
\textbf{Prévision de demande énergétique} | Python, scikit-learn
\begin{itemize}[leftmargin=*]
\item Modélisation de séries temporelles et évaluation des erreurs.
\end{itemize}
\cvsection{Compétences}
\textbf{Data} : Python, SQL, ETL/ELT\\
\textbf{Machine Learning} : scikit-learn, évaluation
\cvsection{Formation}
\textbf{École du numérique -- Diplôme d'ingénieur} \hfill 2021--2026
\cvsection{Langues}
Français natif | Anglais C1
\cvsection{Centres d'intérêt}
Énergie | Lecture
\end{document}
